\subsection{泛映射的实例}

* 以下示例大多将在本系列的其他部分详细讨论。

I. 自由代数结构。设 E 为一个集合，且令 \(\Sigma\) 设代数结构的一个种，由一种或多种合成法则定义。我们取该种的同态作为态射。 \(\Sigma\) 所考虑的，以及作为 \(\alpha\) -映射为从 E 到某个 \(\Sigma\) -集的任意映射（换言之， \(\alpha\{x, s\} = \mathcal{F}(E, x)\) ）。所有常见的代数结构种类都满足 \((\mathrm{CU}_{\mathrm{III}})\) ；除了可除环结构之外，它们还满足 \((\mathrm{CU}_{\mathrm{I}})\) ，以及 \((\mathrm{CU}_{\mathrm{II}})\) 这里是……的一个平凡推论 \((\mathrm{CU}_{\mathrm{I}})\) .

由于一般情况下存在物种的结构 \(\Sigma\) 定义在至少包含两个元素的集合上， \(\alpha\) -映射将E的元素分开，因此E可被视为嵌入到 \(F_{E}\) . \(F_{E}\) 被称为由E生成的自由 \(\Sigma\) -集合。因此在代数中，我们谈论自由半群、自由群、自由模和自由代数。

\begin{enumerate}
\def\labelenumi{\Roman{enumi}.}
\setcounter{enumi}{1}
\item
  分式环与分式域。设E为含单位元的交换环，并设S为E中不包含0的乘法封闭子集。我们取 \(\Sigma\) 设为带单位元的交换环的结构类，态射为将单位元映到单位元的环同态。 \(\alpha\) 映射将为同态 \(\varphi\) ，将E映入带单位元的交换环A，使得 \(\varphi(1)=1\) 并且 \(\varphi(S)\) 仅包含A中的单位。条件 \((\mathbf{Q}\mathbf{M}_{\mathbf{II}}), (\mathbf{C}\mathbf{U}_{\mathbf{I}})\) 通过 \((\mathbf{C}\mathbf{U}_{\mathbf{III}})\) （与 \(\mathfrak{a}=\operatorname{Card}(\mathbf{E})\operatorname{Card}(\mathbf{N}))\) 立即得到验证。因此，泛映射问题总有解。 \((\mathbf{F}_{\mathbf{E}},\varphi_{\mathbf{E}})\) ，但一般而言 \(\varphi_{E}\) 不是单射的。最常出现的情形是E为整环；在这种情况下， \(\varphi_{E}\) 是单射的。此外，如果我们取 \(S = E - \{0\}\) ，那么 \(F_{E}\) 是一个域，称为E的分式域。
\item
  两个模的张量积。设 E 为两个模 A × B 在具有单位元的交换环 C 上的乘积。取 Σ 为 C-模结构的种类，态射为线性映射，α-映射为从 A × B 到某个 C-模的双线性映射。条件 (QM \(_{II}\) ）显然被满足，（CU \(_{I}\) ）到（CU \(_{III}\) ）亦然（其中 a = Card (E) Card (C) Card (N)）。对应于对 (A, B) 的泛 C-模 F \(_{E}\) 称为 A 与 B 的张量积，记作 A ⊗ B。泛映射 φ \(_{E}\) 记作 (x, y) → x ⊗ y；它是双线性的，但一般不是单射。
\item
  模的算子环的扩张。设 A 为带单位元的交换环，B 为 A 的包含 A 的单位元的子环，E 为 B-模。该种 \(\Sigma\) 是 A-模结构的种，态射为 A-线性映射，而 \(\alpha\) -映射为 E 到 A-模的 B-线性映射。对应于 B-模 E 的泛 A-模 \(F_{E}\) 称为将 E 的算子环 B 扩张到 A 而得到的模。
\end{enumerate}

V. 一致空间的完备化。设 E 为一个一致空间。取 \(\Sigma\) 为完备 Hausdorff 一致空间结构的种，态射为一致连续映射，而 \(\alpha\) -映射为 E 到完备 Hausdorff 一致空间的一致连续映射。此处完备 Hausdorff 一致空间的 \(\Sigma\) -容许子集为闭子集（关于空间的拓扑），且条件 \((\mathrm{QM}_{\mathrm{II}})\) 并且 \((\mathrm{CU}_{\mathrm{I}})\) 通过 \((\mathrm{CU}_{\mathrm{III}})\) 被满足（其中 \(a = 2^{2^{\operatorname{Card}(E)}}\) ）。完备 Hausdorff 一致空间（在同构意义下）是与 E 相关联的 Hausdorff 一致空间的完备化（《一般拓扑学》第二章，第3节，第7号）。

\begin{enumerate}
\def\labelenumi{\Roman{enumi}.}
\setcounter{enumi}{5}
\item
  Stone-Čech 紧化。设 E 为一个完全正则空间。 \(\Sigma\) 是紧致空间结构的物种，态射为连续映射（从一个紧致空间到另一个紧致空间），且 \(\alpha\) -映射是将E映入紧空间的连续映射。 \(\Sigma\) -可容许子集再次是闭子集，且条件 \((\mathrm{QM}_{\mathrm{II}})\) , \((\mathrm{CU}_{\mathrm{I}})\) 通过 \((\mathrm{CU}_{\mathrm{III}})\) 容易验证（基数与例V中的相同）。紧空间 \(F_{E}\) 是（在同构意义下）通过关于使E到区间的所有连续映射一致的最粗一致性完备化E而得到的“Stone-Čech紧化”。 {[}0, 1{]} R上的映射是一致连续的（一般拓扑学，第九章，§1，习题7）；映射 \(\varphi_{E}\) 是单射的，因为E中任意两个不同的点都可以通过E到 {[}0, 1{]}.
\item
  自由拓扑群。设E为完全正则空间，设 \(\Sigma\) 为豪斯多夫拓扑群结构的范畴，其态射为连续同态；并取 \(\alpha\) -映射为从 E 到豪斯多夫拓扑群的连续映射。条件（ \(QM_{II}\) ）和（ \(CU_{I}\) ）到（ \(CU_{III}\) 容易验证，
\end{enumerate}

\[
\mathfrak {a} = \text { Card } (\mathbf {E}) \text { Card } (\mathbf {N}).
\]

豪斯多夫拓扑群 \(F_{E}\) 它是E的这个泛映射问题的解，称为由空间E生成的自由拓扑群。由于E的任意两个不同点可以通过E到豪斯多夫拓扑群R的连续映射分离，映射 \(\varphi_{E}\) 是单射；可以证明 \(\varphi_{E}\) 是从 E 到子空间 \(\varphi_{\mathrm{E}}(\mathrm{E})\) 的 \(F_{E}\) (*) 的同胚。我们不取 \(\Sigma\) 作为豪斯多夫拓扑群的结构种，我们也可以采用其他结构种，例如豪斯多夫拓扑阿贝尔群、紧群、豪斯多夫拓扑环、豪斯多夫拓扑向量空间（在拓扑除环上，将其视为辅助基集）等。

\begin{enumerate}
\def\labelenumi{\Roman{enumi}.}
\setcounter{enumi}{7}
\item
  关于拓扑群上的概周期函数。设E为一个拓扑群。取 \(\Sigma\) 作为紧致群结构的范畴，其态射为连续同态，且 \(\alpha\) -映射为从 E 到紧群的连续同态。条件 (QM \(_{II}\) ), (CU \(_{I}\) ）到（CU \(_{III}\) ) 被满足，且 \(a = 2^{2^{\text{Card(E)}}}\) . 紧群 F \(_{E}\) 它是E的这个泛映射问题的解，称为与E相联系的紧群；该映射 \(\varphi_{E}\) 不一定是单射的。E 上每个形如 \(g \circ \varphi_{E}\) ，其中 g 是定义在 \(F_E\)，被称为 E 上的概周期函数。
\item
  阿尔巴尼斯簇。设E是一个代数簇，令 \(\Sigma\) 设为与E同基域的阿贝尔簇结构物种（阿贝尔簇是带有代数群结构的完备代数簇；它必然是交换的）。态射是从一个阿贝尔簇到另一个阿贝尔簇的有理映射（每个态射必然是同态与平移的复合）。 \(\alpha\) -映射是从E到阿贝尔簇的有理映射。条件（CU \(_{I}\) 不满足，然而这个关于 E 的泛映射问题却存在一个解 \(F_{E}\) ，称为 E 的 Albanese 簇。一般地，有理映射 \(\varphi_{E}\) 并非单射。*
\end{enumerate}

